%%%PAGE 114 rot=0 legibility=good summary=极值类习题五道：平抛落到抛物面的最小动能、洛伦兹力下球在碗内圆周运动的最小B、风力作用下的最小风速、斜面上系统牛二律求地面摩擦力极值、斜面上求最小拉力
%%%BLOCK id=114-01 ch=04 sec=平抛 kind=problem alt=06,90 cont=none y=0.03-0.33
% 本题只有标题“交弧用轨迹方程”、图和“求 Ek min”，没有写出完整题干
\begin{problem}
\textbf{交\unclear{弧}用轨迹方程}

%FIG p114_a 0.01 0.085 0.235 0.19
\notefig[0.3]{figures/p114_a.png}

求 $E_{k\min}$
\begin{solution}
\begin{align*}
y&=-\frac{g}{2v_0^2}x^2+2h=\frac{x^2}{2h}\\
\therefore\ y&=\frac{2v_0^2h}{gh+v_0^2}
\end{align*}
先表达 $x^2=\dfrac{4hv_0^2h}{gh+v_0^2}=\dfrac{4v_0^2h^2}{gh+v_0^2}$

\struck{$y=\unclear{\ldots}$}\quad $y=\dfrac{1}{2h}x^2\ \Rightarrow\ y$．
\begin{align*}
E_k&=mg(2h-y)+\frac{1}{2}mv_0^2\qquad\swarrow\\
&=\frac{2mg^2h^2}{gh+v_0^2}+\frac{1}{2}mv_0^2\\
&=\frac{2mg^2h^2}{gh+v_0^2}+\frac{1}{2}m(gh+v_0^2)-\frac{1}{2}mgh\\
&\ge2\sqrt{2mg^2h^2\cdot\frac{1}{2}m}-\frac{1}{2}mgh=\frac{3}{2}mgh
\end{align*}
\end{solution}
\end{problem}
%%%END
%%%BLOCK id=114-02 ch=11 sec=洛伦兹力·圆周运动 kind=problem alt=04 cont=none y=0.25-0.46
\begin{problem}[2]
% 图p114_b：N、qvB、mg 的力三角形（含θ）；图p114_c：球在碗内表面运动的受力图（N、qvB、mg）
%FIG p114_b 0.225 0.248 0.35 0.31
\notefig[0.25]{figures/p114_b.png}

%FIG p114_c 0.01 0.305 0.19 0.40
\notefig[0.3]{figures/p114_c.png}

球能圆周运动 求 $B_{\min}$ 及对应速率
\begin{solution}
【法1】\quad $\dfrac{mv^2}{r}-qvB+mg\tan\theta=0$\qquad $\Delta\ge0\ \Rightarrow\ B$

【法2】\quad $\left\{\begin{array}{l}qvB-N\sin\theta=\dfrac{mv^2}{r}\\[2mm] N\cos\theta=mg\end{array}\right.$

$\therefore\ B=\dfrac{mv}{qr}+\dfrac{mg\tan\theta}{qv}\ge2\sqrt{\dfrac{m^2g\tan\theta}{q^2r}}$ % 此行右端被拍摄边缘截断（可见 m^2 g tan…/q^2 r），按可见部分转录

当 $\dfrac{mv}{qr}=\dfrac{mg\tan\theta}{qv}$ 时取等
\end{solution}
\end{problem}
%%%END
%%%BLOCK id=114-03 ch=02 sec=共点力平衡 kind=problem alt=90 cont=none y=0.44-0.58
\begin{problem}[3]
% 图p114_d内含标注：水平线上的T、θ、mg、F方向等（受力示意）；图p114_e：F、θ、mg、T 的力三角形
%FIG p114_d 0.04 0.44 0.262 0.54
\notefig[0.35]{figures/p114_d.png}

$F=\rho SV^2\sin^2\theta$\quad 求 $V_{\min}$
\begin{solution}
\[mg=F\cos\theta\]
\[\therefore\ V=\sqrt{\frac{mg}{\rho S\sin^2\theta\cos\theta}}\qquad \tan\theta=\frac{\sqrt{2}}{2}\ \text{时}\ V_{\min}\] % CHECK: 若 F∝sin^2θ，使 sin^2θ cosθ 最大应得 tanθ=√2；原稿写作 √2/2（拍照看似 √2 over 2）

%FIG p114_e 0.70 0.495 0.84 0.585
\notefig[0.25]{figures/p114_e.png}
\end{solution}
\end{problem}
%%%END
%%%BLOCK id=114-04 ch=03 sec=连接体与整体隔离 kind=problem alt=90 cont=none y=0.58-0.72
\begin{problem}[4]
\textbf{$\sin2\theta=2\sin\theta\cos\theta$ 及辅助角公式}

% 图p114_f内含标注：g sinθ、μ=0（或M=0，原稿字形难辨）、m、M、θ
%FIG p114_f 0.025 0.645 0.24 0.71
\notefig[0.3]{figures/p114_f.png}

当 $\theta=?$ 时 $f_{\text{地}}$ max\qquad 系统牛二律
\begin{solution}
\[f_{\text{地}}=M\cdot0+ma\cos\theta=mg\sin\theta\cos\theta\qquad \theta=45^\circ\ \text{时取极大值}\]
\end{solution}
\end{problem}
%%%END
%%%BLOCK id=114-05 ch=03 sec=斜面 kind=problem alt=02,90 cont=none y=0.72-0.90
% 原稿序号处“4”被涂改为“5”
\begin{problem}[5]
% 图p114_g内含标注：N、F、f、mg、φ、30°、μ=tan30°
%FIG p114_g 0.085 0.735 0.215 0.84
\notefig[0.3]{figures/p114_g.png}

$\to a=\dfrac{g}{2}$\quad 求最小拉力 $F$
\begin{solution}
\[\left\{\begin{array}{l}F\cos\alpha-mg\sin30^\circ-f=ma\\ N+F\sin\alpha=mg\cos30^\circ\\ f=\mu N\end{array}\right.\qquad F=\frac{mg\sin30^\circ+\mu mg\cos30^\circ+ma}{\sqrt{1+\mu^2}\,\sin(\alpha+\varphi)}\]
\[\tan\varphi=\frac{1}{\mu}\]
当 $\sin(\alpha+\varphi)=1$ 时 $F$ 最小
\end{solution}
\end{problem}
%%%END
%%%PAGEEND 114
